\section{Window Model} 
\label{sec:window_model}
% In order to cluster the evolving data stream in a more efficient way, it is better to only process recent data rather than the whole one. 
% In most cases, more recent information from the stream can reflect the changes in temporal clusters. Thus, this information can be used to detect and explain the evolution of the process under observation, which can improve the algorithm's ability to handle the evolving data stream. 
% To this end, the window model is used to determine the most recent input data for processing. 
The window model can be categorized into three main types as illustrated in Figure~\ref{fig:window_model}. 
% They are landmark (\landmark), sliding (\sliding), and damped (\damped) window models respectively.
For each kind of window model, there are further two types of window implementations: count-based and time-based. 
% Following prior works~\cite{Stream_Clustering_Survey01}, 
We base on the count-based window when implementing \dsc algorithms in our benchmark as it is more frequently used~\cite{Stream_Clustering_Survey01}.


\begin{figure}[t]
\centering
\includegraphics[width=0.7\textwidth]{Figures/Design Aspects/Window Model.pdf}
\caption{Three Types of Window Model.}
\label{fig:window_model}
\end{figure}

\subcompact
\subsection{Landmark Window Model (\landmark)}
\label{subsec:landmark}
% The \landmark handles objects in the form of disjoint portions of the data stream, which is separated by relevant objects such as the time-point or the number of objects observed in the previous window. 
% All objects that arrived after the landmark are considered as part of the current window until the next landmark is reached. When a new landmark is reached, all objects in the previous window are discarded, only those arriving after the new landmark are processed. 
As shown in Figure~\ref{fig:window_model}(a), \landmark collects the data between two landmarks into a window for clustering. When reaching the current landmark, the temporal clustering result will be stored before clearing the current summarizing data structure. After that, the algorithm will continue performing clustering based on the new `blank' structure until reaching the next landmark. 
It is hard to define a proper landmark configuration for \landmark.
% The key challenge of \landmark is that . 
If the two landmarks are spaced very close together, many evolving activities in the stream will not be captured inside the window, leading to poor clustering accuracy. On the contrary, if the two landmarks are set far from each other, a great number of data points will be involved inside the window for computation and the efficiency will decrease significantly.


\subcompact
\subsection{Sliding Window Model (\sliding)}
\label{subsec:sliding}
\sliding is usually implemented as a queue following a First-In-First-Out processing strategy as depicted in Figure~\ref{fig:window_model}(b).
% stores the most recent objects in \change{the window with a predefined window size}
% a data structure, whose size is fixed and defined at the beginning. 
% , where the first object entering the window will be the first one to exit. 
Similar to the \landmark, the \sliding only considers data objects whose timestamp falls within the current window range with older ones discarded. 
In contrast to the \landmark whose window size can be dynamically determined by the landmark, the size of the \sliding is predefined with a fixed value.
% it is usually employed in scenarios when we are only interested in the most recent objects. 
The non-discriminate nature of data insertion and deletion results in a smaller number of clusters by comparison and thus reduces computation time and resources. However, this also means the clustering accuracy will greatly decrease especially under a small window size configuration.
% will be affected by the presence of outliers and \cef in the stream.

% On the contrary, the length of the \sliding needs to be set at the very beginning and will not be changed any longer. Every time when new $k$ points arrive, they will be immediately inserted into the window and the window will also discard its first $k$ points for keeping the length unchanged. 

\subcompact
\subsection{Damped Window Model (\damped)}
\label{subsec:damped}
For both the \landmark and the \sliding, all data objects in the window are associated with equal importance. In contrast, to\margi{R1O1}ensure the more recent data points in the window are always being treated with higher priority, the \damped associate objects in the data stream with weights that decay with a function $f(\alpha, \lambda)$ over time. Meanwhile, unlike the other two types of window models, \damped does not discard objects as time goes on. While taking all data into account may bring in high accuracy, this causes efficiency issues. Furthermore, since the $\alpha$ and $\lambda$ parameters are predefined, the decay function is always kept the same to determine the priority, and the \damped may be more susceptible to some stream characteristics such as \oef.
% Add some discussions of the periodical update process under \damped. 
\begin{comment}
\tony{Notably, some existing algorithms with \damped such as \dstream~\cite{DStream:2007} assume that only one data comes in every time unit. In this case, they only need to update the weight once within a computed time gap, which may greatly improve its efficiency. However, since the gap value is determined by the current number of clusters, this profit may disappear after processing more and more evolving data stream.}
\end{comment}

% In contrast, the \damped associates objects in the data stream with weights that decrease over time, which are usually determined by decay functions. 
% % One typical decay function used for the \damped is the exponential decay function: 
% % \(f(t) = 2^{-\lambda(t_c-t_o)}\), where \(t_c\) is the current time point and \(t_o\) is the creation time. 
% This ensures the most recent information in the data stream is always treated with higher priority. 
% % Different from landmark and sliding window, 
% The \damped does not discard objects as time goes on. Instead, it gradually reduces their significance by associating them with lower weights. While taking all data into account ensures a good level of accuracy, this causes efficiency issues. Furthermore, it is more susceptible to the frequency of evolution and the number of outliers in the data stream.

% Density window model gives every input data a unique weight decaying with time to ensure those most recent data being treated with priority. Landmark and \sliding both similarly collect a bunch of most recent data into a window buffer and use window length and time difference respectively to determine whether to process the collected data using \dsc algorithms.
% In most cases, the recent information from the stream can reflect the changes of temporal clusters. Thus, this information can be used to detect and explain the evolution of the process under observation, which can improve the algorithm's ability to handle the evolving data stream.
% Only when the current window is full will DSC algorithms be executed. 
% Currently, there are three main kinds of window models, they are landmark, sliding and damped window \wangxin{model} respectively.

\textbf{Discussion.}
% Although the three types of window model can all be applied to keep the most recent data from the stream, the different window length makes their clustering behaviour quite different from others. 
Although these window models all try to prioritize the most recent data from streams, their specific strategies are quite different leading to different clustering behaviours. For a better clustering efficiency, \sliding may be the best choice since its fixed window size largely reduces the number of clusters for maintenance. On the contrary, although both \landmark and \damped may not be as efficient as the \sliding, they can bring in more accurate clustering results. Finally, for all three types of window models, their clustering behaviour is sensitive to their specific window configurations, which have not been studied in depth in the literature.
% {First}, the \landmark and \damped have no constraint on their window length, retaining much more information than the \sliding. They essentially bring a tradeoff between accuracy and efficiency. The more information retained, the higher the accuracy, and the lower the efficiency.
% {Second}, it can be a difficult task to set correct window configurations. For the \landmark, it is hard to determine the best landmark. While for the \damped, its decaying function has fixed parameters and can not be adaptive to varying workload characteristics.

% Specifically, as discussed before,\sliding can be helpful to improve the clustering efficiency but may lower the total accuracy. On the contrary, for landmark and damped window model, since they have no constraint on their window length, their clustering accuracy is higher than the sliding one. However, their infinite length of window may also bring in a terrible efficiency problem. Additionally, for \landmark, it is difficult to determine the best landmark. For the damped window model, its usage of decaying function may also damage the original clustering information, which makes itself more susceptible to varying workload characteristics.